\subsection{Charts and atlases. Manifolds.}

5.1.1. Let X be a set. A chart on X is a triple \(c = (U, \varphi, E)\) , where U is a subset of X, E a Banach space and \(\varphi\) a bijection from U onto an open subset of E. We say that U is the domain of the chart c. If E is of finite dimension n, we say that c is of dimension n. Otherwise, we set \(\dim c = +\infty\) .

5.1.2. We say that two charts \(c = (\mathbf{U}, \varphi, \mathbf{E})\) and \(c' = (\mathbf{U}', \varphi', \mathbf{E}')\) of \(X\) are \(C^r\)-compatible (or simply compatible when there can be no ambiguity about \(r\)) when the following conditions are satisfied:

\begin{enumerate}
\def\labelenumi{(\alph{enumi})}
\item
  \(\varphi (\mathbf{U}\cap \mathbf{U}^{\prime})\) (resp. \(\varphi^{\prime}(\mathbf{U}\cap \mathbf{U}^{\prime}))\) is open in \(\mathbf{E}\) (resp. \(\mathbf{E}'\) );
\item
  the map \(\varphi \circ \varphi'^{-1}\) (resp. \(\varphi' \circ \varphi^{-1}\) ) of \(\varphi'(U \cap U')\) onto \(\varphi(U \cap U')\) (resp. of \(\varphi(U \cap U')\) onto \(\varphi'(U \cap U')\) ) is of class \(C^r\) (cf. 2.3.1, 3.2.1 and 4.2.1).
\end{enumerate}

5.1.3. We call a \(\mathbf{C}^{r}\)-atlas (or simply atlas) of a set X a set of charts of X that are pairwise \(\mathbf{C}^{r}\)-compatible and whose domains have union X. Two atlases A and B of X are said to be \(\mathbf{C}^{r}\)-equivalent if \(A \cup B\) is an atlas. The relation of \(\mathbf{C}^{r}\)-equivalence between atlases is an equivalence relation.

5.1.4. Let \(\mathfrak{S}\) be a set of Banach spaces. We say that an atlas \(\mathcal{A}\) is of type \(\mathfrak{S}\) if one has \(E \in \mathfrak{S}\) for every chart \(c = (U, \varphi, E)\) of \(\mathcal{A}\) . Similarly, one says that an atlas \(\mathcal{A}\) is of Hilbert type (resp. of Hilbert type with countable dimension) if \(E\) is a Hilbert space (resp. Hilbert space of countable type) for every chart \((U, \varphi, E)\) of \(\mathcal{A}\) .

5.1.5. One calls a K-manifold of class \(\mathbf{C}^{r}\) (or manifold of class \(\mathbf{C}^{r}\) over K, or simply manifold when there can be no ambiguity about K and r) a set X equipped with a class of equivalent atlases (Sets, chap. II, § 6, no. 9) for the relation of \(\mathbf{C}^{r}\)-equivalence. An atlas of this class is called an atlas of the manifold X. A chart belonging to an atlas of X is called a chart of the manifold X. A chart of X whose domain contains a point a \(\in X\) is called a chart of X at a. A chart centered at a is a chart (U, \(\varphi\), E) of X at a such that \(\varphi(a) = 0\).

If X is a set and A an atlas of X, the set X equipped with the equivalence class of A is called the manifold defined by A.

When \(r \neq \omega\) (hence \(K = R\) ), a manifold of class \(C^r\) is sometimes called a differentiable manifold. A manifold of class \(C^\omega\) is also called an analytic manifold over \(K\) (or also a \(K\) -analytic manifold).

When moreover K = R (resp. C, \(Q_{p}\) ), one also says real analytic manifold (resp. complex analytic manifold, p-adic analytic manifold).

5.1.6. Let X be a manifold. The subsets of X that are unions of domains of charts of X form the set of open sets for a topology on X, called the topology underlying the manifold structure of X. For every chart \(c = (U, \varphi, E)\) of X, the map \(\varphi\) is a homeomorphism from the open set U, equipped with the topology induced by that of X, onto the open set \(\varphi(U)\) of E.

The topological space underlying X is a Baire space. When K is equal to R or C, it is locally connected.

Let X be a manifold and let A be an atlas of X. For the topological space X to be Hausdorff, it is necessary and sufficient that the following condition be satisfied: for any charts (U, \(\varphi\) , E) and (V, \(\psi\) , F) belonging to the atlas A, the graph of the mapping \(\psi \circ \varphi^{-1}\) of \(\varphi(\mathrm{U} \cap \mathrm{V})\) to \(\psi(\mathrm{U} \cap \mathrm{V})\) is closed in \(\varphi(\mathrm{U}) \times \psi(\mathrm{V})\) .

Let X be a manifold; suppose that the topological space X is regular. Then for every point \(a \in X\) there exists a chart (U, \(\varphi\) , E) of X at a possessing the following property: for a subset Y of U to be closed in X, it is necessary and sufficient that its image \(\varphi(Y)\) be closed in E. If the space X is paracompact, there exists on X a distance compatible with the topology of X and making X into a complete metric space.

5.1.7. Let S be a set of Banach spaces. A manifold is said to be of type S (resp. of Hilbertian type, resp. of Hilbertian type of countable dimension) if it possesses an atlas of type S (resp. of Hilbertian type, resp. of Hilbertian type of countable dimension). If S is reduced to a single element E, a manifold of type S is also called a pure manifold of type E.

A pure manifold of dimension n is called a pure manifold of type \(K^{n}\) . A manifold is said to be locally of finite dimension if it is of type S with \(S = \{K^{n}; n \in N\}\) .

5.1.8. Let X be a manifold and let \(x \in X\) . The dimension (finite or \(+\infty\) ) of a chart of X at x depends only on x. It is called the dimension of X at x and is denoted \(\dim_{x}X\) . The dimension of X, denoted \(\dim X\) , is the upper bound of the \(\dim_{x}X\) for \(x \in X\) .

The function \(x \mapsto \dim_x X\) is locally constant. The manifold \(X\) is locally of finite dimension (resp. pure of dimension \(n\) ) if and only if one has \(\dim_x X < +\infty\) (resp. \(\dim_x X = n\) ) for all \(x \in X\) .

5.1.9. Suppose K is locally compact. Let X be a separated manifold. Then, for X to be locally of finite dimension, it is necessary and sufficient that X be locally compact.

5.1.10. Let X be a manifold and let \(\xi^{1},\ldots,\xi^{n}\) be maps from a subset U of X into K. One says that \(\xi=(\xi^{1},\ldots,\xi^{n})\) is a coordinate system of X in U if the triple (U, \(\xi\) , K \(^{n}\) ) is a chart of X; this chart is also denoted (U; \(\xi\) ) or (U; \(\xi^{1},\ldots,\xi^{n}\) ). If \(a\in\mathbf{U}\) , we also say that \(\xi\) is a coordinate system of X at a; if moreover \(\xi^{i}(a)=0\) for \(i=1,2,\ldots,n\) , one says that the coordinate system \(\xi\) is centered at a.

